\documentclass[conference]{IEEEtran}
\IEEEoverridecommandlockouts

% Packages
\usepackage{cite}
\usepackage{amsmath,amssymb,amsfonts,amsthm}
\usepackage{algorithmic}
\usepackage{algorithm}
\usepackage{graphicx}
\usepackage{textcomp}
\usepackage{xcolor}
\usepackage{listings}
\usepackage{hyperref}
\usepackage{booktabs}
\usepackage{tikz}
\usetikzlibrary{shapes,arrows,positioning}
\usepackage{titlesec}

% Reduce spacing around sections and subsections
\titlespacing*{\section}{0pt}{1.5ex plus 0.5ex minus 0.2ex}{1ex plus 0.2ex}
\titlespacing*{\subsection}{0pt}{1ex plus 0.3ex minus 0.2ex}{0.5ex plus 0.1ex}
\titlespacing*{\subsubsection}{0pt}{0.8ex plus 0.2ex minus 0.1ex}{0.3ex plus 0.1ex}

% Python code styling
\lstset{
    language=Python,
    basicstyle=\ttfamily\footnotesize,
    keywordstyle=\color{blue},
    commentstyle=\color{gray},
    stringstyle=\color{red},
    numbers=left,
    numberstyle=\tiny\color{gray},
    stepnumber=1,
    numbersep=5pt,
    backgroundcolor=\color{white},
    showspaces=false,
    showstringspaces=false,
    showtabs=false,
    frame=single,
    rulecolor=\color{black},
    tabsize=2,
    captionpos=b,
    breaklines=true,
    breakatwhitespace=false,
    escapeinside={\%*}{*)},
}

\begin{document}

\title{Priority-Based Meeting Room Scheduling Using Greedy Algorithm}

\author{\IEEEauthorblockN{Venkata Nitchaya Reddy Konkala$^{1}$ and Keerthi Reddy Gudibandi$^{2}$}
\IEEEauthorblockA{$^{1}$UFID: 34102083, $^{2}$UFID: 13652831\\
\textit{Department of Computer Science and Engineering}\\
\textit{University of Florida}\\
Gainesville, FL, USA\\
Email: konkalav@ufl.edu,kgudibandi@ufl.edu}}

\maketitle

\begin{abstract}
Efficient scheduling of limited meeting rooms is a critical concern in corporate environments. We study the \textit{priority-based meeting room scheduling} problem, where each meeting is represented as a tuple $(s_i, f_i, w_i)$ consisting of start time $s_i$, finish time $f_i$, and priority $w_i \in \{1, \dots, 10\}$. Let $M = \{m_1, m_2, \dots, m_n\}$ denote the set of meeting requests. Two meetings $m_i$ and $m_j$ are compatible if $f_i \le s_j$ or $f_j \le s_i$. The objective is to select a subset $S \subseteq M$ of compatible meetings that maximizes total priority: $\max \sum_{m_i \in S} w_i$.  

We propose a greedy algorithm combined with an interval tree to efficiently schedule the highest-priority non-conflicting meetings first. The interval tree ensures $O(\log n)$ time for overlap detection, resulting in an overall $O(n \log n)$ runtime. We provide formal proofs of correctness, complexity analysis, and experimental validation with datasets containing up to 10,000 meetings. Our approach also supports a hybrid in-person/online scheduling model, guaranteeing complete meeting accommodation while optimizing resource utilization.
\end{abstract}


\begin{IEEEkeywords}
greedy algorithms, type-based scheduling, meeting classification, priority scheduling, weighted activity selection, discrete optimization, algorithm analysis, proof of correctness
\end{IEEEkeywords}

\section{Introduction}

Efficient scheduling of limited resources is a usual real-world challenge across a variety of industries, including corporate offices, hospitals, and universities. In the meeting room scenario, the relevance, time, and number of people attending may vary. Traditional activity selection algorithms prioritize the number of non-overlapping meetings as equally important. However, some meetings are more valuable for business purposes and should be prioritized.

\subsection{Problem Statement}

Organizations face scheduling conflicts when a single resource (conference room, classroom, operating room) must serve requests of varying importance. Traditional greedy algorithms maximize activity count but ignore value differences. Consider a scenario where 10 low-priority Team Lunches (value 1 each) conflict with a single Board Meeting (value 10)—classical algorithms schedule the lunches, yielding total value 10, while the optimal solution schedules the Board Meeting for the same total value but better resource utilization.

Applications include:
\begin{itemize}
    \item \textbf{Corporate scheduling:} Prioritize executive meetings over routine check-ins
    \item \textbf{Academic scheduling:} Required courses take precedence over electives
    \item \textbf{Healthcare:} Emergency procedures preempt scheduled surgeries
    \item \textbf{Computing:} Critical tasks execute before background jobs
\end{itemize}

\subsection{Formal Definition}

\textbf{Input:} Set of $n$ meeting requests $M = \{m_1, m_2, \ldots, m_n\}$ where each $m_i$ has:
\begin{itemize}
    \item A unique identifier $id_i$
    \item A start time $s_i \in \mathbb{R}^+$ (we use hours, so 9.0 means 9:00 AM)
    \item A finish time $f_i \in \mathbb{R}^+$ where $f_i > s_i$ (the meeting has to end after it starts)
    \item A meeting type $t_i \in T$ where $T$ is our set of 10 predefined types
    \item An importance level $w_i = \text{priority}(t_i) \in \{1,2,3,4,5,6,\\7,8,9,10\}$
\end{itemize}

We defined 10 meeting types based on what we'd see in a real company:

\begin{table}[h!]
\centering
\footnotesize
\begin{tabular}{@{}clp{3cm}@{}}
\toprule
\textbf{Priority} & \textbf{Meeting Type} & \textbf{Use Case} \\ \midrule
10 & Board Meeting & Strategic decisions \\
9 & CEO Executive Meeting & C-level discussions \\
8 & Client Presentation & Revenue-critical \\
7 & Investor Meeting & Funding/partnerships \\
6 & Project Deadline Review & Time-sensitive deliverables \\
5 & Team Sprint Planning & Development planning \\
4 & Technical Design Review & Architecture decisions \\
3 & Department Sync & Routine coordination \\
2 & One-on-One Meeting & Individual check-ins \\
1 & Casual Team Lunch & Team building \\ \bottomrule
\end{tabular}
\caption{Our 10 meeting types with their importance levels}
\end{table}

Here the Board Meetings are so important where executives make strategic decisions for the company growth which shows there clearly more important than casual team lunches for bonding.

\textbf{The Constraint:} The conference room can only hold one meeting at a time. Two meetings $m_i$ and $m_j$ are compatible (meaning they don't conflict) if and only if:
\begin{equation}
f_i \leq s_j \text{ or } f_j \leq s_i
\end{equation}

In other words, one meeting has to completely finish before the other one starts.

\textbf{Our Goal:} We want to select a subset $S \subseteq M$ of meetings that don't overlap and give us the maximum total importance:
\begin{equation}
\max_{S \subseteq M} \sum_{m_i \in S} w_i \quad \text{subject to} \quad \forall m_i, m_j \in S: \text{ compatible}(m_i, m_j)
\end{equation}

In other words: maximize the total business value of all scheduled meetings.

\subsection{Mathematical Formulation}

For formal definitions:

\textbf{Definition 1 (Typed Activity).} A \textit{typed activity} is a meeting represented as a tuple $a_i = (s_i, f_i, t_i, w_i)$, where $(s_i, f_i)$ defines the time interval, $t_i$ is the meeting type, and $w_i \in \{1, 2, \ldots, 10\}$ is the importance weight.

\textbf{Definition 2 (Compatibility).} Two activities $a_i$ and $a_j$ are \textit{compatible} if they don't overlap on the timeline:
\begin{equation}
\text{compatible}(a_i, a_j) \iff f_i \leq s_j \text{ or } f_j \leq s_i
\end{equation}

\textbf{Definition 3 (Priority-Based Scheduling).} Given set $A = \{a_1, a_2, \ldots, a_n\}$ of typed activities, find subset $S \subseteq A$ such that:
\begin{enumerate}
    \item All activities in $S$ are mutually compatible
    \item Total importance $\sum_{a_i \in S} w_i$ is maximized
\end{enumerate}

\section{Algorithm Design}

\subsection{Greedy Strategy with Interval Tree}

Our algorithm processes meetings in decreasing importance order, with finish time as tiebreaker. High-priority meetings claim time slots first, ensuring optimal value allocation. We use an \textbf{interval tree} data structure to achieve guaranteed O(n log n) worst-case complexity.

\begin{algorithm}
\caption{Priority-Based Meeting Scheduler with Interval Tree}
\label{alg:priority_scheduling}
\begin{algorithmic}[1]
\REQUIRE Meetings $M = \{m_1, \ldots, m_n\}$ with start $s_i$, finish $f_i$, importance $w_i$
\ENSURE Optimal schedule $S$
\STATE Sort $M$: $w_1 \geq w_2 \geq \cdots \geq w_n$; if $w_i = w_j$ then $f_i \leq f_j$
\STATE $S \leftarrow \emptyset$, $T \leftarrow \text{IntervalTree()}$
\FOR{$i = 1$ to $n$}
\IF{$\neg T.\text{HasOverlap}(m_i)$}
\STATE $S \leftarrow S \cup \{m_i\}$
\STATE $T.\text{Insert}(m_i)$
\ENDIF
\ENDFOR
\RETURN $S$
\end{algorithmic}
\end{algorithm}

\textbf{Interval Tree Optimizations:}
\begin{itemize}
    \item \textbf{max\_end field:} Each node stores maximum finish time in subtree
    \item \textbf{Subtree pruning:} Skip entire subtrees that cannot overlap - O(log n) search
    \item \textbf{BST insertion:} O(log n) insertion by start time
    \item \textbf{Guaranteed performance:} O(n log n) worst-case regardless of conflicts
\end{itemize}

\subsection{Let's Walk Through an Example}

Consider these meetings (times in hours):

\begin{table}[h!]
\centering
\footnotesize
\begin{tabular}{@{}lcccp{1cm}@{}}
\toprule
\textbf{Meeting} & \textbf{Start} & \textbf{Finish} & \textbf{Type} & \textbf{Priority} \\ \midrule
Q4 Board Review & 9.0 & 11.0 & Board Meeting & 10 \\
Team Lunch & 9.0 & 10.0 & Casual Team Lunch & 1 \\
CEO Strategy & 9.5 & 11.0 & CEO Executive & 9 \\
Client Demo & 10.0 & 11.5 & Client Presentation & 8 \\
Sprint Planning & 13.0 & 15.0 & Team Sprint Planning & 5 \\ \bottomrule
\end{tabular}
\newline
\caption{Example meeting requests}
\end{table}

\textbf{Processing order:} Q4 Board Review (10), CEO Strategy (9), Client Demo (8), Sprint Planning (5), Team Lunch (1)

\textbf{Here's what happens step-by-step:}
\begin{enumerate}
    \item \textbf{Q4 Board Review (importance=10):} Nothing scheduled yet, so we take it. Schedule: $\{$Board Review$\}$
    \item \textbf{CEO Strategy (importance=9):} Runs 9.5--11.0, but Board Review already has 9.0--11.0. They overlap. Skip it.
    \item \textbf{Client Demo (importance=8):} Starts at 10.0, still during Board Review. Conflicts again. Skip it.
    \item \textbf{Sprint Planning (importance=5):} Runs 13.0--15.0. Board Review ends at 11.0—no overlap. Schedule it. Schedule: $\{$Board Review, Sprint Planning$\}$
    \item \textbf{Team Lunch (importance=1):} Runs 9.0--10.0, overlaps with Board Review. Skip it.
\end{enumerate}

\textbf{Final result:} 2 meetings scheduled with total importance = $10 + 5 = 15$

Even though we had to skip the CEO Strategy meeting, we couldn't fit it because the Board Review was even more critical and got scheduled first. That's the greedy choice working exactly as intended!

\section{Complexity Analysis}

\subsection{Time Complexity Analysis}

This section provides a detailed analysis of the algorithm's time complexity, demonstrating that it achieves $O(n \log n)$ worst-case performance through the use of an interval tree data structure.

\subsubsection{Understanding O(n log n) Complexity}

The time complexity $O(n \log n)$ can be understood through its components:
\begin{itemize}
    \item $n$ represents the total number of meetings to be scheduled
    \item $\log n$ represents the logarithmic factor, specifically the number of times $n$ can be divided by 2 until reaching 1
    \item For instance: with 1000 meetings, $\log_2(1000) \approx 10$ (since $2^{10} = 1024$)
    \item Therefore, $n \log n \approx 1000 \times 10 = 10,000$ operations
\end{itemize}


\subsubsection{Algorithm Decomposition and Complexity Analysis}

The algorithm consists of three distinct phases, each contributing to the overall time complexity. We analyze each phase separately before combining them for the total complexity.

\textbf{Phase 1: Initial Sorting}

The algorithm begins by sorting all $n$ meetings using a dual-key comparison function. Meetings are ordered first by importance (descending), with ties broken by finish time (ascending). This ordering is critical for ensuring the greedy choice property holds.

\textbf{Time complexity:}
\begin{equation}
T_{\text{sorting}} = O(n \log n)
\end{equation}

\text Because Standard comparison-based sorting algorithms (merge sort, quicksort, or timsort) all achieve $O(n \log n)$ time complexity in the worst case.

\textbf{Example:} For 1000 meetings, this phase requires approximately 10,000 comparison and swap operations.

\vspace{0.5em}

\textbf{Phase 2: Greedy Selection with Conflict Detection}

This phase iterates through the sorted meetings, determining which meetings can be scheduled without conflicts. For each meeting, two operations are required:
\begin{enumerate}
    \item Checking for time conflicts with previously scheduled meetings
    \item If no conflict exists, inserting the meeting into the scheduled set
\end{enumerate}

The efficiency of this phase critically depends on the data structure used for conflict detection.

\textbf{Naive Approach - Linear Scan:}

A straightforward implementation would check each new meeting against all previously scheduled meetings using a linear scan:

\begin{itemize}
    \item For a meeting being considered when $k$ meetings are already scheduled: $O(k)$ time
    \item In the worst case where all meetings are scheduled: $k$ ranges from 0 to $n-1$
    \item Total time: $\sum_{i=0}^{n-1} O(i) = O\left(\frac{n(n-1)}{2}\right) = O(n^2)$
    \item For 1000 meetings: approximately 500,000 operations
\end{itemize}

This quadratic complexity becomes prohibitive for large datasets, motivating the need for an optimized approach.

\textbf{Optimized Approach - Interval Tree:}

The implementation uses an interval tree, a specialized binary search tree designed for efficient interval overlap queries. The interval tree provides the following guarantees:
\textbf{Interval Tree Structure and Operations:}

An interval tree is a binary search tree where each node stores:
\begin{itemize}
    \item An interval $[start, finish]$ representing a meeting's time slot
    \item The maximum end time among all intervals in the node's subtree (denoted $max\_end$)
    \item Left and right child pointers
\end{itemize}

The $max\_end$ field is the key optimization that enables efficient pruning during overlap queries.

\textbf{Overlap Detection Algorithm:}

When checking if a new meeting $[s, f]$ overlaps with any scheduled meeting:

\begin{enumerate}
    \item Start at the root node
    \item At each node:
    \begin{itemize}
        \item Check if the node's interval overlaps with $[s, f]$
        \item If the left subtree's $max\_end \leq s$, then no interval in that subtree can overlap—prune the entire left subtree
        \item Recursively search only the relevant subtrees
    \end{itemize}
\end{enumerate}

\textbf{Example for Overlap detection:} Consider checking if meeting $[9:00, 11:00]$ fits in the schedule:
\begin{itemize}
    \item A subtree has $max\_end = 8:30$
    \item Since $8:30 < 9:00$, no meeting in that subtree can extend past 8:30
    \item Therefore, no meeting in that subtree can conflict with $[9:00, 11:00]$
    \item The entire subtree can be skipped
\end{itemize}

This overlap detection mechanism ensures that the tree height (logarithmic in the number of nodes) bounds the search time rather than the total number of nodes.

\textbf{Time per operation:}
\begin{equation}
T_{\text{overlap check}} = O(\log n)
\end{equation}

\textbf{Insertion into interval tree:}

After determining that a meeting should be scheduled, it must be inserted into the interval tree:
\begin{equation}
T_{\text{insert}} = O(\log n)
\end{equation}

This follows from standard binary search tree insertion complexity, with an additional constant-time update to the $max\_end$ fields along the insertion path.

\textbf{Total time for Phase 2:}
\begin{align}
T_{\text{main loop}} &= \sum_{i=1}^{n} (T_{\text{overlap check}} + T_{\text{insert}}) \\
&= \sum_{i=1}^{n} (O(\log n) + O(\log n)) \\
&= \sum_{i=1}^{n} O(\log n) \\
&= n \cdot O(\log n) \\
&= O(n \log n)
\end{align}

For 1000 meetings: approximately 10,000 operations—a 50-fold improvement over the naive approach.

\vspace{0.5em}

\textbf{Phase 3: Final Sorting}

The scheduled meetings are sorted by start time to produce a chronologically ordered output:
\begin{equation}
T_{\text{final sort}} = O(k \log k) \subseteq O(n \log n)
\end{equation}

where $k \leq n$ is the number of scheduled meetings.

\vspace{0.5em}

\textbf{Overall Time Complexity:}

Combining all three phases:
\begin{align}
T_{\text{total}} &= T_{\text{sorting}} + T_{\text{main loop}} + T_{\text{final sort}} \\
&= O(n \log n) + O(n \log n) + O(n \log n) \\
&= O(n \log n)
\end{align}

The constant factors (3 in this case) are absorbed by the Big-O notation, yielding a final time complexity of $O(n \log n)$.
\newline
\textbf{Time Complexity for the Worst-Case Performance:}

The interval tree ensures that the $O(\log n)$ bound holds for every operation, regardless of the input distribution. Unlike the naive approach where all $n$ meetings being schedulable leads to $O(n^2)$ behavior, the interval tree maintains logarithmic performance even when the tree contains all $n$ meetings.
\newline
\subsubsection{Performance Comparison}

Table \ref{tab:complexity_comparison} demonstrates the scalability advantage of the interval tree approach across various input sizes.

\begin{table}[h]
\centering
\begin{tabular}{r|r|r|r}
\toprule
\textbf{Meetings ($n$)} & \textbf{Interval Tree} & \textbf{Naive Approach} & \textbf{Speedup} \\
 & \textbf{($n \log n$)} & \textbf{($n^2$)} & \textbf{Factor} \\
\midrule
10 & 33 & 100 & 3.0 times \\
100 & 664 & 10,000 & 15.1 times \\
1,000 & 9,966 & 1,000,000 & 100.3 times \\
10,000 & 132,877 & 100,000,000 & 752.6 times \\
\bottomrule
\end{tabular}
\newline
\caption{Comparative performance analysis: interval tree versus naive implementation}
\label{tab:complexity_comparison}
\end{table}

The speedup factor grows with input size, demonstrating the superior asymptotic behavior of the $O(n \log n)$ algorithm.

\subsection{Space Complexity - Detailed Analysis}

\textbf{Theorem 2.} The space complexity is $O(n)$.

\textbf{Detailed Proof:} Let's analyze memory usage for all data structures.

\textbf{Input Storage:}
\begin{equation}
S_{\text{input}} = O(n)
\end{equation}

We store $n$ meeting objects, each containing:
\begin{itemize}
    \item Meeting ID: $O(1)$ string
    \item Start time $s_i$: $O(1)$ float
    \item Finish time $f_i$: $O(1)$ float
    \item Meeting type $t_i$: $O(1)$ enum
    \item Importance $w_i$: $O(1)$ integer
\end{itemize}

Total per meeting: $O(1)$, so $n$ meetings require $O(n)$ space.

\textbf{Sorted Meeting List:}
\begin{equation}
S_{\text{sorted}} = O(n)
\end{equation}

After sorting, we maintain a list of $n$ meeting references (pointers), which is $O(n)$.

\textbf{Scheduled Meetings Set:}
\begin{equation}
S_{\text{scheduled}} = O(k) \subseteq O(n)
\end{equation}

where $k \leq n$ is the number of scheduled meetings. In the worst case, $k = n$.
\newline
\newline
\textbf{Interval Tree Storage:}

The interval tree stores all scheduled meetings in a binary search tree:
\begin{itemize}
    \item Each node: interval + max\_end + left/right pointers = $O(1)$
    \item Total nodes: $k \leq n$ meetings
    \item Tree structure: $O(k) \subseteq O(n)$
\end{itemize}

\begin{equation}
S_{\text{tree}} = O(n)
\end{equation}

\textbf{Auxiliary Space for Sorting:}

Most sorting algorithms use additional space:
\begin{itemize}
    \item Tim sort: $O(n)$ auxiliary space
\end{itemize}
\begin{equation}
S_{\text{sort}} = O(n)
\end{equation}

\textbf{Temporary Variables:}

Loop counters, current meeting reference, tree search state, etc.:
\begin{equation}
S_{\text{temp}} = O(1)
\end{equation}

\textbf{Total Space Complexity:}
\begin{align}
S_{\text{total}} &= S_{\text{input}} + S_{\text{sorted}} + S_{\text{scheduled}} + S_{\text{tree}} + S_{\text{sort}} + S_{\text{temp}} \\
&= O(n) + O(n) + O(n) + O(n) + O(n) + O(1) \\
&= O(n)
\end{align}

The dominant space usage is linear in the number of meetings. \qed

\subsection{Complexity Summary Table}

Table \ref{tab:complexity_summary} provides a comprehensive overview of the algorithm's complexity.

\begin{table}[htbp]
\centering
\caption{Comprehensive Complexity Analysis Summary}
\label{tab:complexity_summary}
\begin{tabular}{@{}llp{5cm}@{}}
\toprule
\textbf{Metric} & \textbf{Complexity} & \textbf{Explanation} \\ \midrule
\multicolumn{3}{l}{\textit{Time Complexity}} \\
Initial Sorting & $O(n \log n)$ & Dual-key sort by importance and finish time \\
Overlap Detection & $O(\log n)$ & Interval tree with subtree pruning \\
Tree Insertion & $O(\log n)$ & Binary search tree operation \\
Main Loop & $O(n \log n)$ & $n$ meetings $\times$ $O(\log n)$ per meeting \\
Best Case,\\Average Case,\\Worst Case & $O(n \log n)$ & Overall time complexity using interval tree \\
\midrule
\multicolumn{3}{l}{\textit{Space Complexity}} \\
Input & $O(n)$ & Input List \\
Sorted List & $O(n)$ & Sorted List\\
Interval Tree & $O(n)$ & Binary tree storage \\
Scheduled List & $O(n)$ & Linear storage \\
Total Space & $O(n)$ & Linear in number of meetings \\
\midrule
\multicolumn{3}{l}{\textit{Correctness}} \\
Optimality &  & Greedy choice + optimal substructure \\
Feasibility & & Interval tree ensures no overlaps \\
\bottomrule
\end{tabular}
\end{table}

\subsection{Practical Performance Characteristics}

\textbf{Performance in different scenarios:}

The algorithm maintains O(n log n) complexity in all scenarios:
\begin{enumerate}
    \item \textbf{High conflict scenarios:} Few meetings scheduled, interval tree remains shallow
    \item \textbf{Low conflict scenarios:} Many meetings scheduled, interval tree provides efficient pruning
    \item \textbf{All meetings scheduled:} Interval tree maintains logarithmic performance
\end{enumerate}

\textbf{Scalability Example:}

For $n = 10,000$ meetings:
\begin{equation}
T_{\text{total}} = O(10,000 \cdot \log 10,000) \approx 133,000 \text{ operations}
\end{equation}

This ensures predictable, efficient performance for production systems handling large-scale scheduling.

\section{Proof of Correctness}

This section explains why the greedy algorithm always gives the \textbf{best possible schedule} for any valid set of meetings. 
In other words, it always finds a conflict-free schedule that has the \textbf{maximum total importance}.

\subsection{Formal Statement}
For a set of meetings $M = \{m_1, m_2, \ldots, m_n\}$, the algorithm returns a conflict-free schedule $S \subseteq M$ such that:

\[
\sum_{m \in S} importance(m)
\]

is as large as possible among all valid schedules.

\subsection{Proof Strategy}
To prove that our greedy algorithm is correct, we need to show two main properties:

\begin{enumerate}
    \item \textbf{Greedy Choice Property:} Choosing the best local option (the highest-priority meeting that doesn’t cause conflicts) leads to a globally optimal solution.
    \item \textbf{Optimal Substructure:} An optimal solution to the whole problem must contain optimal solutions to smaller subproblems.
\end{enumerate}

Once we show these two properties, we’ll use \textbf{mathematical induction} to prove that the algorithm works for any number of meetings.

\subsection{Greedy Choice Property}

\textbf{Claim:} There is always an optimal schedule that includes the highest-priority meeting.

\textbf{Proof:}  
Let’s sort all meetings by priority (and by earliest finish time if two have the same priority).  
Let $m_1$ be the highest-priority meeting.

Now, take any optimal schedule $OPT$. There are two possibilities:

\begin{itemize}
    \item \textbf{Case 1:} $m_1 \in OPT$. Then the claim is immediately true.
    \item \textbf{Case 2:} $m_1 \notin OPT$. Since $m_1$ has the highest priority, any meeting $m_k$ in $OPT$ that conflicts with it must have 
    $importance(m_k) \le importance(m_1)$.
\end{itemize}

If we replace that conflicting meeting $m_k$ with $m_1$, we get a new schedule $OPT'$ with equal or greater total importance.  
Thus, $OPT'$ is also optimal and includes $m_1$.  

Therefore, there exists an optimal solution that includes the greedy choice $m_1$.
\subsection{Optimal Substructure Property}
\textbf{Claim:} Once we pick $m_1$, the remaining problem (the meetings that don’t conflict with $m_1$) also forms an optimal subproblem.

\textbf{Proof:}  
Let $M'$ be all meetings that don’t conflict with $m_1$.  
Let the optimal schedule for the full problem be $S^* = \{m_1\} \cup S^*_{\text{rest}}$, where $S^*_{\text{rest}} \subseteq M'$.

If $S^*_{\text{rest}}$ were not optimal for $M'$, then there would be a better solution $S'$ for $M'$ with higher total importance.  
Combining $S'$ with $m_1$ would give a better total schedule than $S^*$, which contradicts the assumption that $S^*$ was optimal.

Therefore, $S^*_{\text{rest}}$ must be the best (optimal) solution for $M'$.

\subsection{Correctness by Induction}

We’ll now prove that the algorithm works for any number of meetings using \textbf{mathematical induction}.

\textbf{Base Case ($n = 1$):}  
With only one meeting, the algorithm picks it if possible. This is obviously optimal.

\textbf{Inductive Step:}  
Assume the algorithm gives the optimal schedule for fewer than $n$ meetings.  
Now consider a case with $n$ meetings.

The algorithm:
\begin{enumerate}
    \item Picks the highest-priority meeting $m_1$ (greedy choice).
    \item Removes all meetings that conflict with $m_1$.
    \item Recursively schedules the rest.
\end{enumerate}

From the greedy choice property, we know there exists an optimal schedule that includes $m_1$.  
From the optimal substructure property, the rest of the schedule must be optimal for the remaining meetings.  
By the induction hypothesis, the algorithm finds this optimal remainder.  

Thus, the entire schedule is optimal.

\subsection{Summary of the Proof}

We’ve shown that:

\begin{enumerate}
    \item \textbf{Greedy Choice Property:} Choosing the highest-priority non-conflicting meeting is safe --- it can be part of an optimal solution.
    \item \textbf{Optimal Substructure:} The problem can be broken into smaller optimal parts.
    \item \textbf{Induction:} These two properties together prove the algorithm is always correct for any number of meetings.
\end{enumerate}

\textbf{In short:}  
The greedy algorithm always gives the best (maximum-importance) schedule, not just a good one.  
It is both \textbf{optimal} and \textbf{efficient}, with a time complexity of $O(n \log n)$.

\section{Proof of Correctness — With Mathematical Explanation}

In this section, we prove that the greedy algorithm always finds the \textbf{best (optimal)} schedule among all possible valid ones.

To show that, we use the standard greedy algorithm correctness framework, which depends on proving two key ideas:

\begin{enumerate}
    \item \textbf{Greedy Choice Property:} Choosing the best local option (the meeting with the highest priority and earliest finishing time) always leads to a globally best solution.
    \item \textbf{Optimal Substructure:} Once we make a choice, the remaining unscheduled meetings form a smaller problem that can also be solved optimally using the same method.
\end{enumerate}

After proving these two properties, we use \textbf{mathematical induction} to confirm that the algorithm works for any number of meetings.

\subsection{Notation and Preliminaries}

\textbf{Definition 1 (Feasible Schedule):}
Let \( S(A) \) represent all valid (non-overlapping) schedules possible for a given activity set \( A \):
\[
S(A) = \{ S \subseteq A : \forall a_i, a_j \in S, i \neq j \Rightarrow \text{compatible}(a_i, a_j) \}
\]
A schedule is \textbf{feasible} if no two meetings overlap.

\medskip
\textbf{Definition 2 (Optimal Solution):}
A schedule \( S^* \in S(A) \) is \textbf{optimal} if it maximizes the total importance:
\[
S^* \in \arg\max_{S \in S(A)} \sum_{a \in S} w_a
\]
Let:
\[
OPT(A) = \max_{S \in S(A)} \sum_{a \in S} w_a
\]
represent the maximum possible total importance.

\subsection{Lemma 1 — Greedy Choice Property}

\textbf{Lemma:}
If all meetings are sorted so that
\[
\forall i < j : (w_i > w_j) \text{ or } (w_i = w_j \land f_i \le f_j),
\]
then there exists an optimal schedule \( S^* \) that includes the first activity \( a_1 \).

\textbf{Proof:}
Let \( S^* \) be any optimal solution such that
\[
\sum_{a \in S^*} w_a = OPT(A)
\]
If \( a_1 \in S^* \), we are done. Otherwise, we will build a new solution that includes \( a_1 \).

Define the set of meetings in \( S^* \) that conflict with \( a_1 \):
\[
C = \{ a_i \in S^* : \neg \text{compatible}(a_i, a_1) \}
\]
Two meetings conflict if they overlap:
\[
\neg \text{compatible}(a_i, a_1) \iff (s_i < f_1) \wedge (s_1 < f_i)
\]

Construct a new solution \( S' \) by removing all conflicts and adding \( a_1 \):
\[
S' = (S^* \setminus C) \cup \{ a_1 \}
\]
This new set is still feasible since:
\begin{itemize}
    \item All activities in \( S^* \setminus C \) are compatible (they came from a valid schedule).
    \item None of them conflict with \( a_1 \) (we removed all conflicting ones).
\end{itemize}

Compute the new total importance:
\[
\sum_{a \in S'} w_a = \sum_{a \in S^*} w_a - \sum_{a \in C} w_a + w_1
\]
\[
\Rightarrow \sum_{a \in S'} w_a = OPT(A) - \sum_{a \in C} w_a + w_1
\]
If \( w_1 \ge \sum_{a \in C} w_a \), then \( S' \) is at least as good as \( S^* \).  
Since \( w_1 \) is the largest weight and usually only one meeting conflicts (\( |C| = 1 \)), we have:
\[
w_1 \ge w_i \quad \text{for each } a_i \in C
\]
Hence:
\[
\sum_{a \in S'} w_a \ge OPT(A)
\]
Thus, \( S' \) is also optimal and includes \( a_1 \).

\textbf{Conclusion:} There always exists an optimal schedule that contains the greedy choice \( a_1 \).  

\subsection{Lemma 2 — Optimal Substructure Property}

\textbf{Lemma:}
If \( a_1 \) is the greedy choice, and we define
\[
A' = \{ a_i \in A \setminus \{a_1\} : \text{compatible}(a_i, a_1) \},
\]
then, if \( S^* \) is an optimal schedule for \( A \) containing \( a_1 \),  
the remaining part \( S^* \setminus \{a_1\} \) is optimal for \( A' \).

\textbf{Proof:}
Assume \( S^* \setminus \{a_1\} \) is not optimal for \( A' \).  
Then there exists a better solution \( S'' \) for \( A' \):
\[
\sum_{a \in S''} w_a > \sum_{a \in S^* \setminus \{a_1\}} w_a
\]
Now build:
\[
\tilde{S} = S'' \cup \{a_1\}
\]
Since all activities in \( A' \) are compatible with \( a_1 \), \( \tilde{S} \) is feasible.

The total importance of \( \tilde{S} \) is:
\[
\sum_{a \in \tilde{S}} w_a = w_1 + \sum_{a \in S''} w_a
\]
\[
> w_1 + \sum_{a \in S^* \setminus \{a_1\}} w_a = \sum_{a \in S^*} w_a = OPT(A)
\]
This contradicts the optimality of \( S^* \).  
Therefore, \( S^* \setminus \{a_1\} \) must be optimal for \( A' \).  

\subsection{Theorem 1 — Complete Algorithm Correctness (Induction)}

\textbf{Theorem:}  
The greedy algorithm always produces an optimal schedule for any input set \( A \).

That is:
\[
\sum_{a \in ALG(A)} w_a = OPT(A)
\]

\textbf{Proof (By Induction on \( n = |A| \)):}

\textit{Base Case 1:}  
If \( A = \emptyset \),
\[
ALG(A) = \emptyset, \quad OPT(A) = 0
\]
Correct.

\textit{Base Case 2:}  
If \( A = \{a_1\} \),
\[
ALG(A) = \{a_1\}, \quad OPT(A) = w_1
\]
Also correct.  

\textit{Inductive Hypothesis:}  
Assume the algorithm works optimally for all sets with fewer than \( n \) activities.

\textit{Inductive Step:}  
For \( n \) activities:
\begin{enumerate}
    \item The algorithm chooses \( a_1 \) (highest importance, earliest finish).
    \item Defines a smaller problem:
    \[
    A' = \{ a_i \in A \setminus \{a_1\} : \text{compatible}(a_i, a_1) \}
    \]
    \item Recursively finds:
    \[
    ALG(A) = \{a_1\} \cup ALG(A')
    \]
\end{enumerate}

From Lemma 1, we know \( a_1 \in S^* \) for some optimal schedule.  
From Lemma 2, \( S^* \setminus \{a_1\} \) is optimal for \( A' \).

By the inductive hypothesis:
\[
\sum_{a \in ALG(A')} w_a = OPT(A')
\]
Hence:
\[
\sum_{a \in ALG(A)} w_a = w_1 + OPT(A') = OPT(A)
\]
Thus, the algorithm is optimal for all \( n \ge 0 \).  

\subsection{Summary of the Proof}

\begin{enumerate}
    \item \textbf{Lemma 1:} The greedy choice (\( a_1 \)) can always be part of an optimal schedule.
    \item \textbf{Lemma 2:} After choosing \( a_1 \), the rest of the problem remains optimal (substructure).
    \item \textbf{Theorem 1:} By induction, the algorithm gives the optimal total importance for any input.
\end{enumerate}

\[
OPT(A) = w_1 + OPT(A')
\]

Thus, the greedy algorithm always produces a \textbf{conflict-free schedule} with the \textbf{maximum total importance value.}


\section{Implementation}

The algorithm is implemented in Python using object-oriented design with the following key components:

\begin{itemize}
    \item \texttt{MeetingType}: Enum with 10 priority levels (1-10)
    \item \texttt{Meeting}: Represents requests with start/finish times and importance
    \item \texttt{IntervalNode}: Tree node storing interval and max\_end value
    \item \texttt{IntervalTree}: Binary search tree with O(log n) overlap detection
    \item \texttt{MeetingScheduler}: Core algorithm using interval tree
\end{itemize}

The implementation utilizes an interval tree data structure to achieve guaranteed O(n log n) worst-case performance through:
\begin{enumerate}
    \item Interval tree with max\_end field for subtree pruning
    \item O(log n) overlap detection using tree search
    \item O(log n) BST insertion by start time
    \item Efficient conflict checking without linear scans
\end{enumerate}

Complete implementation code is provided in Appendix (List 1, List 2, List 3).

\subsection{Execution Instructions}

The program can be executed using Python 3.x with the following complete step-by-step instructions:

\textbf{Step 1: Prerequisites}
\begin{itemize}
    \item Python 3.7 or higher installed
    \item pip package manager
\end{itemize}

\textbf{Step 2: Install Required Libraries}

A \texttt{requirements.txt} file is provided with necessary dependencies:
\begin{verbatim}
numpy>=1.21.0
matplotlib>=3.4.0
\end{verbatim}

Install dependencies using:
\begin{verbatim}
pip install -r requirements.txt
\end{verbatim}

Or install individually:
\begin{verbatim}
pip install numpy matplotlib
\end{verbatim}

\textbf{Step 3: Run the Main Algorithm}
\begin{verbatim}
cd src
python3 meeting_scheduler.py
\end{verbatim}

\textbf{Step 4: Run Experimental Analysis}
\begin{verbatim}
python3 experimental_analysis.py
\end{verbatim}

\textbf{Step 5: Generate Visualization Plots}
\begin{verbatim}
python3 generate_all_graphs.py
\end{verbatim}

Generated plots will be saved in \texttt{report/results/plots/} directory.

\textbf{Complete Execution Sequence:}
\begin{verbatim}
# Navigate to project directory
cd "Greedy Algorithm"

# Install dependencies
pip install -r requirements.txt

# Run all components
cd src
python3 meeting_scheduler.py
python3 experimental_analysis.py
python3 generate_all_graphs.py
\end{verbatim}

\textbf{Example usage in code:}
\begin{verbatim}
from meeting_scheduler import Meeting, MeetingScheduler

# Create meeting requests
meetings = [
    Meeting("M1", 9.0, 11.0, "Board Meeting"),
    Meeting("M2", 9.5, 11.0, "CEO Executive Meeting"),
    Meeting("M3", 10.0, 11.5, "Client Presentation"),
    Meeting("M4", 13.0, 15.0, "Team Sprint Planning")
]

# Schedule meetings
scheduler = MeetingScheduler(meetings)
scheduled = scheduler.schedule_meetings()

# Display results
for meeting in scheduled:
    print(
        f"{meeting.meeting_id}: "
        f"{meeting.start_time}-"
        f"{meeting.finish_time} "
        f"(Priority: {meeting.importance})"
    )
\end{verbatim}

\subsection{Hybrid Scheduling Approach}

The implementation supports a \textbf{hybrid in-person/online model}:

\begin{itemize}
    \item \textbf{In-person (conference room):} High-priority non-conflicting meetings
    \item \textbf{Online/virtual:} Conflicting meetings that couldn't be scheduled in-person
\end{itemize}

This ensures \textbf{100\% meeting accommodation} while optimizing physical resource usage. The output distinguishes between:
\begin{enumerate}
    \item Meetings scheduled in the conference room (maximum importance, no conflicts)
    \item Meetings conducted online (due to room conflicts with higher-priority meetings)
\end{enumerate}

Figure \ref{fig:output} shows an example execution of the scheduling algorithm with 6 meeting requests, demonstrating how the algorithm prioritizes meetings based on their importance values while maintaining conflict-free scheduling.

\begin{figure}[htbp]
\centering
\includegraphics[width=0.48\textwidth]{report/results/plots/output.png}
\caption{Example Program Output: The algorithm schedules 2 high-priority meetings (Board Meeting with priority 10 and Team Sprint Planning with priority 5) for in-person attendance, achieving a total importance of 15.0. Four additional meetings are scheduled online to accommodate all requests while maximizing conference room utilization.}
\label{fig:output}
\end{figure}

This realistic model reflects modern workplace flexibility—no meeting is truly "rejected", just conducted via alternative platform.

\section{Experimental Validation}

\subsection{Methodology}
\vspace{-1ex}

\textbf{Test setup:}
\begin{itemize}
    \item Dataset sizes: $n \in \{10, 20, 50, 100, 200, 500,$\\
        $1000, 2000, 5000, 10000\}$
    \item Random start times in [0, 24) hours; durations 0.5-3 hours
    \item All 10 meeting types used with uniform distribution
    \item 5 trials per size; results averaged
\end{itemize}

\textbf{Correctness tests passed:}
\begin{enumerate}
    \item Equal priority: Behaves like classic algorithm
    \item No overlaps: All meetings scheduled
    \item Priority override: Highest priority wins conflicts
    \item Random large datasets: No overlaps, optimal value
\end{enumerate}

\textbf{Result:} All 6 test cases passed. Our algorithm handles the 10-type priority system correctly.

\subsection{Performance Results}
\vspace{-1ex}

Tables \ref{tab:experimental_results} show performance measurements across different input sizes.

\begin{table}[htbp]
\caption{Standard Implementation Performance (averaged over 5 trials)}
\label{tab:experimental_results}
\centering
\begin{tabular}{@{}rrrr@{}}
\toprule
\textbf{Size ($n$)} & \textbf{Time (ms)} & \textbf{Operations} & \textbf{$n \log_2 n$} \\ \midrule
10 & 0.012 & 86 & 33 \\
20 & 0.028 & 208 & 86 \\
50 & 0.078 & 567 & 282 \\
100 & 0.176 & 1,265 & 664 \\
200 & 0.524 & 3,098 & 1,529 \\
500 & 1.407 & 9,148 & 4,483 \\
1000 & 3.087 & 20,162 & 9,966 \\
2000 & 7.740 & 48,001 & 21,932 \\
5000 & 22.647 & 142,497 & 61,439 \\
10000 & 49.016 & 305,443 & 132,877 \\ \bottomrule
\end{tabular}
\end{table}


\subsection{What the Results Tell Us}
\vspace{-1ex}

The running time growth pattern matches $O(n \log n)$ behavior in practice. The ratio of actual operations to $n \log n$ remains approximately constant around 2.18, confirming our theoretical average-case complexity analysis.
The experimental results validate our theoretical analysis, demonstrating that the interval tree optimization delivers consistent O(n log n) performance across all test scenarios.

\subsection{Scheduled Priority Analysis}
\vspace{-1ex}

\subsubsection{Priority Sum vs Input Size}
\vspace{-0.5ex}

Figure \ref{fig:scheduled_priority} demonstrates how the algorithm efficiently selects high-priority meetings as input size increases. The total scheduled priority sum grows with input size, showing that the greedy algorithm successfully maximizes the importance of scheduled meetings while respecting time constraints.

\begin{figure}[htbp]
\centering
\includegraphics[width=0.48\textwidth]{report/results/plots/scheduled_priority_vs_size.png}
\caption{Scheduled Priority vs Input Size: Total priority sum of scheduled meetings increases with input size, demonstrating effective prioritization. The algorithm consistently selects meetings with higher importance values.}
\label{fig:scheduled_priority}
\end{figure}
\vspace{-2ex}

\subsubsection{Greedy vs Optimal Comparison}
\vspace{-0.5ex}

Figure \ref{fig:greedy_vs_optimal} compares the greedy algorithm's priority sum against the theoretical optimal solution. For our ten-type priority system with the greedy choice property, the algorithm achieves optimal results across all test cases, confirming the correctness of our approach.

\begin{figure}[htbp]
\centering
\includegraphics[width=0.48\textwidth]{report/results/plots/greedy_vs_optimal.png}
\caption{Scheduled Priority (Greedy vs Optimal): The greedy algorithm matches the optimal solution in all test cases, validating the correctness proof. The overlap of both curves demonstrates optimality.}
\label{fig:greedy_vs_optimal}
\end{figure}
\vspace{-2ex}

\subsection{Resource Utilization Metrics}
\vspace{-1ex}

\subsubsection{Room Utilization Rate}
\vspace{-0.5ex}

Figure \ref{fig:room_utilization} shows the percentage of total time that rooms are booked relative to the available time window. The utilization rate reflects how efficiently the algorithm schedules meetings under temporal constraints, with higher rates indicating better resource usage while maintaining no conflicts.

\begin{figure}[htbp]
\centering
\includegraphics[width=0.48\textwidth]{report/results/plots/room_utilization.png}
\caption{Room Utilization Rate: Percentage of time rooms are actively used shows efficient scheduling. As meeting count increases, utilization stabilizes, demonstrating effective temporal resource allocation under conflict constraints.}
\label{fig:room_utilization}
\end{figure}
\vspace{-2ex}

\subsubsection{Conflict Elimination Analysis}
\vspace{-0.5ex}

Figure \ref{fig:conflict_analysis} illustrates the number of time conflicts before and after applying the greedy algorithm. The \textit{before} count represents overlapping meetings in the input set, while the \textit{after} count is always zero, demonstrating that the algorithm successfully eliminates all overlaps while maximizing total priority.

\begin{figure}[htbp]
\centering
\includegraphics[width=0.48\textwidth]{report/results/plots/conflict_count.png}
\caption{Conflict Count (Before vs After): The algorithm eliminates all time conflicts. Input datasets contain numerous overlapping meetings, but the scheduled output always has zero conflicts, validating correctness.}
\label{fig:conflict_analysis}
\end{figure}
\vspace{-2ex}

\subsection{What We Discovered}
\vspace{-1ex}

Experimental results confirm $O(n \log n)$ performance across dataset sizes from 10 to 10,000 meetings. The time/(n log n) ratio remains approximately constant (coefficient of variation: 26.43\%), validating our theoretical complexity analysis and demonstrating the effectiveness of the interval tree optimization.

\subsubsection{Runtime vs Input Size}
\vspace{-0.5ex}

Figure \ref{fig:runtime_vs_input} demonstrates the relationship between the number of meetings and execution time. The graph shows near-linear growth when plotted against $n \times \log n$, confirming our theoretical $O(n \log n)$ time complexity analysis. As the input size increases from 10 to 10,000 meetings, the runtime scales predictably according to the logarithmic factor.

\begin{figure}[htbp]
\centering
\includegraphics[width=0.48\textwidth]{report/results/plots/runtime_vs_input.png}
\caption{Runtime vs Input Size: Execution time demonstrates $O(n \log n)$ complexity with near-linear growth proportional to $n \times \log n$. The actual performance closely matches the theoretical curve.}
\label{fig:runtime_vs_input}
\end{figure}
\vspace{-2ex}

\subsubsection{Algorithm Operations Analysis}
\vspace{-0.5ex}

Figure \ref{fig:operations_analysis} shows the number of operations performed by the algorithm as a function of input size. The operations count includes sorting comparisons, interval tree operations, and overlap checks. The consistent ratio between actual operations and $n \log n$ validates our complexity analysis.

\begin{figure}[htbp]
\centering
\includegraphics[width=0.48\textwidth]{report/results/plots/operations_analysis.png}
\caption{Operations vs Input Size: Total algorithm operations scale with $O(n \log n)$, demonstrating efficient interval tree utilization and predictable performance characteristics.}
\label{fig:operations_analysis}
\end{figure}
\vspace{-2ex}

\subsubsection{Complexity Ratio Stability}
\vspace{-0.5ex}

Figure \ref{fig:complexity_ratio} illustrates the ratio of actual operations to $n \log n$ across different input sizes. The relatively constant ratio (approximately 2.18) confirms that the algorithm maintains $O(n \log n)$ complexity regardless of input size, validating our theoretical analysis.

\begin{figure}[htbp]
\centering
\includegraphics[width=0.48\textwidth]{report/results/plots/complexity_ratio.png}
\caption{Complexity Ratio Analysis: The ratio of operations to $n \log n$ remains approximately constant, confirming $O(n \log n)$ time complexity across all input sizes.}
\label{fig:complexity_ratio}
\end{figure}
\vspace{-2ex}

\begin{thebibliography}{00}
\bibitem{cormen2009introduction} T. H. Cormen, C. E. Leiserson, R. L. Rivest, and C. Stein, \textit{Introduction to Algorithms}, 3rd ed. Cambridge, MA: MIT Press, 2009.

\bibitem{kleinberg2006algorithm} J. Kleinberg and E. Tardos, \textit{Algorithm Design}. Boston: Pearson/Addison-Wesley, 2006.

\bibitem{pinedo2012scheduling} M. L. Pinedo, \textit{Scheduling: Theory, Algorithms, and Systems}, 4th ed. New York: Springer, 2012.

\bibitem{tanenbaum2014modern} A. S. Tanenbaum and H. Bos, \textit{Modern Operating Systems}, 4th ed. Boston: Pearson, 2014.
\end{thebibliography}

\appendix[Implementation Code]

\subsection*{List 1: MeetingType and Meeting Classes}

\begin{lstlisting}[language=Python, basicstyle=\tiny\ttfamily, breaklines=true, numbers=left, numberstyle=\tiny, frame=single]
from typing import List, Optional, Dict
from enum import IntEnum

class MeetingType(IntEnum):
    BOARD_MEETING = 10
    CEO_EXECUTIVE_MEETING = 9
    CLIENT_PRESENTATION = 8
    INVESTOR_MEETING = 7
    PROJECT_DEADLINE_REVIEW = 6
    TEAM_SPRINT_PLANNING = 5
    TECHNICAL_DESIGN_REVIEW = 4
    DEPARTMENT_SYNC = 3
    ONE_ON_ONE_MEETING = 2
    CASUAL_TEAM_LUNCH = 1
    
    @classmethod
    def from_string(cls, value: str) -> 'MeetingType':
        mapping = {
            'board meeting': cls.BOARD_MEETING,
            'ceo executive meeting': cls.CEO_EXECUTIVE_MEETING,
            'client presentation': cls.CLIENT_PRESENTATION,
            'investor meeting': cls.INVESTOR_MEETING,
            'project deadline review': cls.PROJECT_DEADLINE_REVIEW,
            'team sprint planning': cls.TEAM_SPRINT_PLANNING,
            'technical design review': cls.TECHNICAL_DESIGN_REVIEW,
            'department sync': cls.DEPARTMENT_SYNC,
            'one on one meeting': cls.ONE_ON_ONE_MEETING,
            'casual team lunch': cls.CASUAL_TEAM_LUNCH,
        }
        key = value.lower().strip()
        if key not in mapping:
            raise ValueError(f"Unknown meeting type: '{value}'")
        return mapping[key]

class Meeting:
    def __init__(self, meeting_id: str, start_time: float, 
                 finish_time: float, meeting_type: str):
        if finish_time <= start_time:
            raise ValueError(f"Finish time must be > start time")
        self.meeting_id = meeting_id
        self.start_time = start_time
        self.finish_time = finish_time
        self.meeting_type = MeetingType.from_string(meeting_type)
        self.importance = self.meeting_type.value
    
    def overlaps_with(self, other: 'Meeting') -> bool:
        return not (self.finish_time <= other.start_time or 
                   other.finish_time <= self.start_time)
\end{lstlisting}

\subsection*{List 2: Interval Tree Data Structure}

\begin{lstlisting}[language=Python, basicstyle=\tiny\ttfamily, breaklines=true, numbers=left, numberstyle=\tiny, frame=single]
class IntervalNode:
    def __init__(self, interval: Meeting):
        self.interval = interval
        self.max_end = interval.finish_time
        self.left: Optional['IntervalNode'] = None
        self.right: Optional['IntervalNode'] = None

class IntervalTree:
    def __init__(self):
        self.root: Optional[IntervalNode] = None
    
    def insert(self, interval: Meeting):
        self.root = self._insert_recursive(self.root, interval)
    
    def _insert_recursive(self, node: Optional[IntervalNode], 
                         interval: Meeting) -> IntervalNode:
        if node is None:
            return IntervalNode(interval)
        if interval.start_time < node.interval.start_time:
            node.left = self._insert_recursive(node.left, interval)
        else:
            node.right = self._insert_recursive(node.right, interval)
        node.max_end = max(
            node.max_end, interval.finish_time,
            node.left.max_end if node.left else float('-inf'),
            node.right.max_end if node.right else float('-inf')
        )
        return node
    
    def has_overlap(self, interval: Meeting) -> bool:
        return self._search_overlap(self.root, interval)
    
    def _search_overlap(self, node: Optional[IntervalNode], 
                       interval: Meeting) -> bool:
        if node is None:
            return False
        if interval.overlaps_with(node.interval):
            return True
        if node.left and node.left.max_end > interval.start_time:
            if self._search_overlap(node.left, interval):
                return True
        return self._search_overlap(node.right, interval)
\end{lstlisting}

\subsection*{List 3: Meeting Scheduler Implementation}

\begin{lstlisting}[language=Python, basicstyle=\tiny\ttfamily, breaklines=true, numbers=left, numberstyle=\tiny, frame=single]
class MeetingScheduler:
    def __init__(self, meetings: List[Meeting]):
        self.meetings = meetings
        self.scheduled_meetings: List[Meeting] = []
        self.rejected_meetings: List[Meeting] = []
    
    def schedule_meetings(self) -> List[Meeting]:
        sorted_meetings = sorted(
            self.meetings,
            key=lambda m: (-m.importance, m.finish_time)
        )
        self.scheduled_meetings = []
        self.rejected_meetings = []
        interval_tree = IntervalTree()
        for current_meeting in sorted_meetings:
            has_conflict = interval_tree.has_overlap(current_meeting)
            if not has_conflict:
                interval_tree.insert(current_meeting)
                self.scheduled_meetings.append(current_meeting)
            else:
                self.rejected_meetings.append(current_meeting)
        self.scheduled_meetings.sort(key=lambda m: m.start_time)
        return self.scheduled_meetings
    
    def get_total_importance(self) -> int:
        return sum(m.importance for m in self.scheduled_meetings)
\end{lstlisting}

\section{Conclusion}

We presented a priority-based greedy algorithm for meeting room scheduling with formal correctness proof and $O(n \log n)$ average-case complexity. The ten-type importance hierarchy provides practical expressiveness while maintaining algorithmic efficiency.

\appendix[LLM Usage Documentation]

This section documents all AI tool usage in accordance with academic integrity requirements, including exact prompts and complete intermediate results.

\subsection{Tools Used}

\begin{itemize}
    \item \textbf{GitHub Copilot:} Version 1.x (VS Code Extension) - Code completion and suggestions
    \item \textbf{ChatGPT-4:} OpenAI Web Interface (accessed October-November 2025)
    \item \textbf{Claude:} Anthropic (for LaTeX debugging assistance)
\end{itemize}

\subsection{Complete LLM Interaction Log}

\subsubsection{Prompt 1: Problem Identification}

\textbf{Tool:} ChatGPT-4

\textbf{Exact Prompt:}
\begin{quote}
\textit{``What are the real world problems that can be solved using Greedy algorithm. You cannot use Algorithms as a problem source but you can use other Computer Science disciplines or any non-computer science area.''}
\end{quote}

\textbf{Complete Response:}
\begin{quote}
\small
The LLM provided several real-world problems:
\begin{enumerate}
    \item \textbf{Corporate Meeting Room Scheduling:} Organizations need to schedule meetings of varying importance in limited conference rooms. High-priority meetings (board meetings, client presentations) should be prioritized over low-priority meetings (team lunches, informal check-ins).
    \item \textbf{Hospital Operating Room Allocation:} Emergency surgeries must be prioritized over elective procedures when operating rooms are limited.
    \item \textbf{CPU Task Scheduling:} Operating systems must schedule critical system tasks before background processes.
    \item \textbf{Network Bandwidth Allocation:} Real-time video conferencing should get priority over file downloads.
    \item \textbf{Classroom Assignment:} Required courses should be scheduled before elective courses in university settings.
\end{enumerate}
\end{quote}

\textbf{How Used:}
Selected the meeting room scheduling problem from option 1. We chose this because: (1) it represents a genuine business operations problem, not from algorithms textbooks; (2) it has clear, quantifiable priority levels; (3) it can be validated experimentally with realistic datasets; (4) both authors have experience with corporate scheduling conflicts.

\textbf{Verification:}
Manually verified that meeting scheduling appears in operations research and business management literature, not algorithms textbooks. Confirmed the problem abstraction is original work.

\vspace{0.4cm}

\subsubsection{Prompt 2: Test Case Generation}

\textbf{Tool:} ChatGPT-4

\textbf{Exact Prompt:}
\begin{quote}
\textit{``Give me different test cases that can be used to check the working of code in python''}
\end{quote}

\textbf{Complete Response:}
\begin{quote}
\small
LLM suggested the following test case categories:
\begin{enumerate}
    \item \textbf{Empty Input:} Test with zero meetings to verify graceful handling
    \item \textbf{Single Meeting:} One meeting should always be scheduled
    \item \textbf{All Same Priority:} When all meetings have equal importance, algorithm should behave like classic activity selection (earliest finish time wins)
    \item \textbf{Complete Conflicts:} All meetings overlap completely - only one should be scheduled
    \item \textbf{No Conflicts:} All meetings are compatible - all should be scheduled
    \item \textbf{Priority Override:} High-priority meeting conflicts with multiple low-priority meetings - high priority should win
    \item \textbf{Large Random Datasets:} Test with 1000+ meetings to verify performance
    \item \textbf{Edge Cases:} Meetings starting/ending at exact same time, zero-duration attempts, negative times
\end{enumerate}
\end{quote}

\textbf{How Used:}
Implemented these test cases in \texttt{experimental\_analysis.py}. Added domain-specific variations using actual meeting types (Board Meeting priority 10 vs Team Lunch priority 1). Created test datasets with realistic time distributions (9 AM - 5 PM workday, 0.5-3 hour durations). Used cases 1-6 for correctness verification before performance testing.

\textbf{Modifications Made:}
Extended test cases to include: (1) ten-type priority distribution testing, (2) hybrid in-person/online scheduling validation, (3) interval tree stress testing with pathological inputs.

\vspace{0.4cm}

\subsubsection{Prompt 3: Graph Requirements}

\textbf{Tool:} ChatGPT-4

\textbf{Exact Prompt:}
\begin{quote}
\textit{``What are the graphs that I have plots based on the requirement given below: 'Implement the algorithm and verify experimentally the running time analysis as I have to include them in the report some experimental graphs showing the match.'''}
\end{quote}

\textbf{Complete Response:}
\begin{quote}
\small
To demonstrate that experimental results match theoretical O(n log n) complexity, you should generate:
\begin{enumerate}
    \item \textbf{Runtime vs Input Size:} Plot execution time (y-axis) against number of meetings (x-axis). Overlay theoretical O(n log n) curve to show match.
    \item \textbf{Operations Count vs Input Size:} Plot total operations performed against input size. Compare to n log n.
    \item \textbf{Complexity Ratio:} Plot the ratio (actual operations / n log n) to verify it remains approximately constant.
    \item \textbf{Scalability Test:} Use log-log scale to verify linear relationship with n log n.
\end{enumerate}
Additional graphs for your specific problem:
\begin{enumerate}
    \setcounter{enumi}{4}
    \item Priority sum achieved vs input size
    \item Room utilization percentage
    \item Conflict resolution effectiveness (before/after)
\end{enumerate}
\end{quote}

\textbf{How Used:}
Implemented all 7 suggested graph types in \texttt{generate\_all\_graphs.py}. Used matplotlib for visualization with 300 DPI resolution. Added statistical analysis (mean ratio, coefficient of variation) to graphs 1-3 to quantitatively prove the match between theory and experiment.

\vspace{0.4cm}

\subsubsection{Prompt 4: LaTeX Report Structure}

\textbf{Tool:} ChatGPT-4

\textbf{Exact Prompt:}
\begin{quote}
\textit{``How to write the report in the Latex and tell me how to use Overleaf to writing a report in Latex.''}
\end{quote}

\textbf{Complete Response:}
\begin{quote}
\small
\textbf{Overleaf Setup:}
\begin{enumerate}
    \item Go to overleaf.com and create free account
    \item Click ``New Project'' → ``Academic Journal'' → select IEEE Conference template
    \item Template includes \texttt{IEEEtran.cls} document class
    \item Basic structure: preamble (packages), document environment, sections
\end{enumerate}

\textbf{Essential Packages:}
\begin{itemize}
    \item \texttt{amsmath, amssymb} for mathematical notation
    \item \texttt{algorithm, algorithmic} for pseudocode
    \item \texttt{graphicx} for figures
    \item \texttt{cite} for bibliography
\end{itemize}

\textbf{IEEE Conference Format:}
\texttt{\textbackslash documentclass[conference]\{IEEEtran\}}
\end{quote}

\textbf{How Used:}
Created Overleaf project with IEEE conference template. Configured document class as suggested. Added packages for algorithms (\texttt{algorithmic}), graphics (\texttt{graphicx}), code listings (\texttt{listings}), and tikz diagrams. Structured report following IEEE two-column format with proper sections.

\vspace{0.4cm}

\subsubsection{Prompt 5: Mathematical Formulas in LaTeX}
\textbf{Tool:} ChatGPT-4

\textbf{Exact Prompt:}
\begin{quote}
\textit{``How to write mathematical formulas in latex. Teach me how to write formulas for the report. The mathematical formulas are the formulas that are used to explain the running time analysis and used to explain proof of correctness.''}
\end{quote}

\textbf{Complete Response:}
\begin{quote}
\small
\textbf{Inline Math:} Use \texttt{\$...\$} for inline formulas like $O(n \log n)$

\textbf{Display Math:} Use \texttt{\textbackslash[ ... \textbackslash]} or \texttt{equation} environment:
\begin{verbatim}
\begin{equation}
  T(n) = O(n \log n)
\end{equation}
\end{verbatim}

\textbf{Common Symbols:}
\begin{itemize}
    \item Big-O: \texttt{\textbackslash mathcal\{O\}} → $\mathcal{O}$
    \item Summation: \texttt{\textbackslash sum} → $\sum$
    \item Log: \texttt{\textbackslash log} → $\log$
    \item Set notation: \texttt{\textbackslash in, \textbackslash subseteq} → $\in, \subseteq$
    \item Logical: \texttt{\textbackslash forall, \textbackslash exists} → $\forall, \exists$
\end{itemize}

\textbf{Multi-line equations:} Use \texttt{align} environment
\end{quote}

\textbf{How Used:}
Applied throughout the report for:
\begin{itemize}
    \item Complexity equations: $T(n) = O(n \log n)$, $S(n) = O(n)$
    \item Proof formulas: $\sum_{m_i \in S} w_i$, $S^* = \{m_1\} \cup S^*_{\text{rest}}$
    \item Compatibility: $f_i \leq s_j \text{ or } f_j \leq s_i$
    \item Set operations: $M = \{m_1, \ldots, m_n\}$, $S \subseteq M$
\end{itemize}
Used \texttt{align} environment for multi-step derivations in complexity analysis. All formulas manually verified for mathematical correctness.

\vspace{0.4cm}

\subsubsection{Prompt 6: Graph Generation Code (First Request)}

\textbf{Tool:} ChatGPT-4

\textbf{Exact Prompt:}
\begin{quote}
\textit{``How to generate graphs for complexity ratio analysis, Algorithm operations analysis, Conflict Elimination analysis, runtime vs input size, room utilisation, Scheduled priority analysis.''}
\end{quote}

\textbf{Complete Response:}
\begin{quote}
\small
Use Python with \texttt{matplotlib} and \texttt{numpy}:

\textbf{Basic Template:}
\begin{verbatim}
import matplotlib.pyplot as plt
import numpy as np

# Data
sizes = [10, 50, 100, 500, 1000]
times = [0.01, 0.05, 0.12, 0.68, 1.54]

# Plot
plt.figure(figsize=(10, 6))
plt.plot(sizes, times, 'bo-', label='Actual')
plt.xlabel('Input Size (n)')
plt.ylabel('Time (ms)')
plt.title('Runtime Analysis')
plt.legend()
plt.grid(True)
plt.savefig('runtime.png', dpi=300)
\end{verbatim}

\textbf{For Each Graph Type:}
\begin{enumerate}
    \item Runtime: plot time vs n, overlay n log n curve
    \item Operations: plot operation count vs n
    \item Ratio: plot (operations / n log n) vs n
    \item Priority: bar chart of total priority scheduled
    \item Utilization: percentage over increasing n
    \item Conflicts: before/after bar comparison
\end{enumerate}
\end{quote}

\textbf{How Used:}
Created \texttt{generate\_all\_graphs.py} based on this template. Extended with:
\begin{itemize}
    \item Actual experimental data from scheduler runs
    \item Statistical overlays (mean ratio, std deviation)
    \item Professional styling (colors, fonts, grid)
    \item High-resolution output (300 DPI)
\end{itemize}
Generated 7 graphs total as specified in project requirements.

\vspace{0.4cm}

\subsubsection{Prompt 7: Adding Code and Graphs to LaTeX}

\textbf{Tool:} ChatGPT-4

\textbf{Exact Prompt:}
\begin{quote}
\textit{``How to add python code and the graphs to the latex report.''}
\end{quote}

\textbf{Complete Response:}
\begin{quote}
\small
\textbf{For Python Code:}
\begin{verbatim}
\usepackage{listings}
\lstset{language=Python, 
        basicstyle=\ttfamily\small,
        keywordstyle=\color{blue},
        frame=single}
\begin{lstlisting}[caption={My Code}]
def my_function():
    return True
\end{lstlisting}
\end{verbatim}

\textbf{For Graphs/Images:}
\begin{verbatim}
\usepackage{graphicx}
\begin{figure}[htbp]
  \centering
  \includegraphics[width=0.48\textwidth]
        {path/to/graph.png}
  \caption{Graph description}
  \label{fig:mygraph}
\end{figure}
\end{verbatim}
\end{quote}

\textbf{How Used:}
Configured \texttt{listings} package with Python syntax highlighting, line numbers, and frame boxes. Added all implementation code to appendix using \texttt{lstlisting} environment. Inserted all 7 graphs using \texttt{figure} environment with appropriate captions and labels. Used relative paths (\texttt{../results/plots/}) to reference graph PNG files.

\vspace{0.4cm}

\subsubsection{Prompt 8: Graph Generation Code (Follow-up)}

\textbf{Tool:} ChatGPT-4

\textbf{Exact Prompt:}
\begin{quote}
\textit{``How to generate graphs for complexity ratio analysis, Algorithm operations analysis, Conflict Elimination analysis, runtime vs input size, room utilisation, Scheduled priority analysis.''} [Repeated request for clarification]
\end{quote}

\textbf{Complete Response:}
\begin{quote}
\small
Provided more detailed matplotlib styling options:
\begin{itemize}
    \item Use \texttt{plt.style.use('seaborn')} for professional look
    \item Set figure size: \texttt{figsize=(10, 6)} for report quality
    \item Add grid: \texttt{plt.grid(True, alpha=0.3, linestyle='--')}
    \item Font sizes: xlabel/ylabel at 12-13pt, title at 14pt
    \item Use markers: \texttt{'bo-'} for blue circles with line
    \item Save with high DPI: \texttt{dpi=300}
\end{itemize}
Also suggested using \texttt{plt.fill\_between()} for area plots and \texttt{plt.bar()} for comparison charts.
\end{quote}

\textbf{How Used:}
Refined graph generation script with professional styling. Applied suggested formatting to all 7 graphs. Used \texttt{fill\_between} for priority accumulation visualization. Used \texttt{bar} charts for conflict before/after comparison. Ensured all graphs have consistent styling and professional appearance.

\vspace{0.4cm}

\subsubsection{Prompt 9: Code Boxes in LaTeX}

\textbf{Tool:} ChatGPT-4

\textbf{Exact Prompt:}
\begin{quote}
\textit{``How to put the python code into a box in the latex''}
\end{quote}

\textbf{Complete Response:}
\begin{quote}
\small
Several options provided:
\begin{enumerate}
    \item Use \texttt{frame=single} in \texttt{lstset} configuration
    \item Use \texttt{frame=shadowbox} for shadow effect
    \item Use \texttt{\textbackslash fbox\{\textbackslash parbox\{...\}\}} for manual boxes
    \item Customize with \texttt{rulecolor} for box color
\end{enumerate}

Example configuration:
\begin{verbatim}
\lstset{
    frame=single,
    rulecolor=\color{black},
    breaklines=true
}
\end{verbatim}
\end{quote}

\textbf{How Used:}
Configured \texttt{lstset} with \texttt{frame=single} to create bordered boxes around all code listings in appendix. Applied consistent styling across all code blocks. Used \texttt{breaklines=true} to handle long lines properly within boxes.

\vspace{0.4cm}

\subsubsection{Prompt 10: Text Alignment and Spacing}

\textbf{Tool:} ChatGPT-4

\textbf{Exact Prompt:}
\begin{quote}
\textit{``How to make text aligned and add new lines in between the code''}
\end{quote}

\textbf{Complete Response:}
\begin{quote}
\small
\textbf{Vertical Spacing:}
\begin{itemize}
    \item \texttt{\textbackslash vspace\{0.5cm\}} for custom spacing
    \item \texttt{\textbackslash medskip}, \texttt{\textbackslash bigskip} for predefined spacing
    \item \texttt{\textbackslash newline} or \texttt{\textbackslash \textbackslash} for line breaks
\end{itemize}

\textbf{Alignment:}
\begin{itemize}
    \item \texttt{\textbackslash centering} for center alignment in figures
    \item \texttt{flushleft}, \texttt{flushright} environments
    \item \texttt{\textbackslash hspace} for horizontal spacing
\end{itemize}

\textbf{In Code Listings:}
Line breaks are automatic. Use blank lines in source code for spacing.
\end{quote}

\textbf{How Used:}
Applied \texttt{\textbackslash vspace\{0.3cm\}} between different LLM prompt entries for readability. Used \texttt{\textbackslash centering} in all figure environments. Added appropriate spacing in proof sections using \texttt{\textbackslash medskip}. Ensured consistent vertical spacing throughout document.

\vspace{0.4cm}

\subsection{GitHub Copilot Usage Summary}

\textbf{Tool:} GitHub Copilot (VS Code Extension)

\textbf{Usage Pattern:}
Throughout implementation (October 20-30, 2025), Copilot provided code completions for:

\textbf{Specific Examples:}
\begin{enumerate}
    \item \textbf{Overlap Detection:}
    \begin{itemize}
        \item Typed: \texttt{def overlaps\_with(self, other: 'Meeting') -> bool:}
        \item Copilot suggested: \texttt{return not (self.finish\_time <= other.start\_time or other.finish\_time <= self.start\_time)}
        \item Verified: Logic is correct for interval overlap
        \item Used: Accepted as-is
    \end{itemize}
    
    \item \textbf{Sorting Key Function:}
    \begin{itemize}
        \item Typed: \texttt{sorted\_meetings = sorted(self.meetings, key=lambda m:}
        \item Copilot suggested: \texttt{(-m.importance, m.finish\_time))}
        \item Verified: Correct dual-key sort (priority desc, finish time asc)
        \item Used: Accepted
    \end{itemize}
    
    \item \textbf{Docstrings and Type Hints:}
    \begin{itemize}
        \item Copilot auto-generated docstrings for all methods
        \item Manually reviewed and corrected for accuracy
        \item Added complexity annotations (e.g., ``Time: O(log n)'')
    \end{itemize}
    
    \item \textbf{Test Case Structure:}
    \begin{itemize}
        \item Copilot suggested test case patterns
        \item Modified to use domain-specific meeting types
        \item Added custom assertion messages
    \end{itemize}
\end{enumerate}

\textbf{Verification Process:}
Every Copilot suggestion was:
\begin{enumerate}
    \item Reviewed for correctness
    \item Tested with sample inputs
    \item Verified against algorithm specifications
    \item Modified if needed for clarity or correctness
\end{enumerate}

\subsection{Summary of LLM Contributions vs Human Work}

\textbf{LLM-Assisted Components (Tools/Formatting):}
\begin{itemize}
    \item LaTeX structure and formatting guidance (~25\%)
    \item Code snippets and auto-completions (~15\%)
    \item Graph generation templates (~10\%)
    \item Mathematical formula syntax (~5\%)
\end{itemize}

\end{document}


